
\documentclass[10pt,a4paper]{article}

\usepackage[
  a4paper,
  left=0.85in,
  right=0.85in,
  top=0.75in,
  bottom=0.85in
]{geometry}

\usepackage[hidelinks]{hyperref}
\usepackage[round,authoryear]{natbib}

\usepackage{indentfirst} % also indents first paragraph after section titles
\setlength{\parindent}{1.5em}
\setlength{\parskip}{0pt}
\setlength{\abovedisplayskip}{4pt}
\setlength{\belowdisplayskip}{4pt}
\setlength{\arraycolsep}{3pt}
\renewcommand{\arraystretch}{1.08}


\begin{document}
\fontsize{9.2}{10.8}\selectfont

\begin{center}
{\Large Thesis Proposal: Functional Forecasting of Turkish Electricity Load Curves}\\
\vspace{0.5em}
{\normalsize Zeynep Bozoklu}\\
\vspace{0.25em}
{\small \href{https://github.com/zebozoklu/turkey-electricity-fda}{GitHub repository}}
\end{center}

\vspace{-0.5em}

\section*{Motivation}

This project studies day-ahead forecasting of Turkish electricity demand, measured hourly in MWh. The raw
quantity of interest is not a single daily number, but the full 24-hour load profile:
how demand evolves overnight, rises in the morning, peaks during the day, and falls
again in the evening. Treating each day as a curve preserves this intraday shape.

The practical question is whether a functional time-series model can improve on the
available official day-ahead forecast. The report therefore compares each proposed
model with two reference points: a simple weekly benchmark, \(Y_{d-7}\), and the
official forecast, \(\widehat Y^{off}_d\). The weekly benchmark checks whether the
model beats a very transparent rule based on the strong weekly rhythm of electricity
use. The official forecast is the main benchmark because it represents the forecast
that would be available in practice.

The proposed baseline contribution is a parsimonious functional forecasting model:
daily curves are represented by FPCA scores, the scores are forecast with lagged
functional dynamics and predetermined calendar, holiday, load-shape, and derivative
features, and forecasts are evaluated in a rolling-origin design. A second-stage
residual adjustment is reported separately as a possible extension to discuss rather
than as the main specification.

\section*{Object and functional representation}

Let \(Y_d(h)\) denote Turkish electricity load on day \(d\), hour \(h=0,\ldots,23\).  
The empirical object is the daily load curve
\[
Y_d=\{Y_d(h):h=0,\ldots,23\}.
\]

The cleaned data are stored as daily curve matrices
\[
Y=[Y_d(h)]_{d,h}, \qquad
\widehat Y^{off}=[\widehat Y^{off}_d(h)]_{d,h}, \qquad
E=[Y_d(h)-\widehat Y^{off}_d(h)]_{d,h},
\]
with 24 hourly observations per retained day.

Each curve is represented by a B-spline basis,
\[
Y_d(h)\simeq \sum_{m=1}^{M}c_{dm}B_m(h),
\qquad M=12.
\]

FPCA gives
\[
Y_d(h)=\mu(h)+\sum_{k=1}^{K}\xi_{dk}\phi_k(h)+e_d(h),
\qquad K=2.
\]

The choice \(K=2\) is based on the empirical FPCA variance decomposition of the
load curves. The first two load components explain \(97.29\%\) of total functional
variation, exceeding the usual 95\% cumulative-variance rule while keeping the
score vector parsimonious. The first component mainly captures the overall daily
level and the second adds the dominant intraday shape contrast. Components three
and four explain only \(1.11\%\) and \(0.99\%\) of additional load variation,
respectively, so they are not retained in the forecasting model. The official
forecast-error FPCA is reported only as a diagnostic description of error shape; it
is not used to choose \(K\).

\begin{center}
\fontsize{8.4}{9.6}\selectfont
\begin{tabular}{rrrrr}
\hline
FPC & Load var. & Load cum. & Error var. & Error cum.\\
\hline
1 & 93.40\% & 93.40\% & 76.60\% & 76.60\%\\
2 & 3.89\% & 97.29\% & 12.00\% & 88.60\%\\
3 & 1.11\% & 98.40\% & 4.62\% & 93.22\%\\
4 & 0.99\% & 99.39\% & 2.88\% & 96.09\%\\
\hline
\end{tabular}
\end{center}

\fontsize{9.2}{10.8}\selectfont

Let
\[
\xi_d=(\xi_{d1},\ldots,\xi_{dK})',
\qquad
\Phi(h)=(\phi_1(h),\ldots,\phi_K(h))'.
\]

Then
\[
Y_d(h)=\mu(h)+\Phi(h)'\xi_d+e_d(h),
\qquad
\widehat Y_d(h)=\widehat\mu(h)+\widehat{\Phi}(h)'\widehat\xi_d.
\]

All forecasts are rolling-origin forecasts: for target day \(d\), estimation uses only information dated \(<d\). The evaluation period starts on 2023-01-01 and contains \(26{,}304\) hourly observations.

\section*{Proposed modeling strategy}

The model sequence is intentionally incremental. Each specification asks whether one
additional source of information explains a remaining part of the forecast error.

First, the daily curve is reduced to a small number of FPCA scores. This is useful
because the 24 hourly observations are highly dependent: morning, afternoon, and
evening demand move together in recurring ways. FPCA summarizes the main level and
shape movements of the whole curve, so the forecasting problem becomes forecasting a
few daily scores rather than 24 separate hourly loads.

The first FPCA models use lagged scores. A weekly lag is natural because electricity
load has a strong day-of-week structure: Mondays tend to resemble previous Mondays,
weekends resemble previous weekends, and so on. Adding a one-day lag allows the model
to react to recent changes in demand level or shape that are not fully captured by
the same weekday one week earlier. This motivates the FPCA VAR(1,7) score model.

The next additions are calendar, holiday, and load-shape predictors. These variables
are included because Turkish electricity demand changes around weekends, national
holidays, religious holidays, bridge days, and seasonal school or summer
periods. Lagged load-shape variables, such as the previous day's mean, peak, range,
and morning ramp, are included because they describe the current demand regime before
the target day is forecast.

Derivative features are then added to describe how quickly the load curve is changing
within the day. Two days can have similar average demand but different ramping
behavior. The derivative summaries and targeted morning, evening, and end-of-day
ramps are meant to capture these slope and transition patterns.

\section*{Forecasting specifications}

Benchmarks:
\[
\widehat Y_d^{SN}(h)=Y_{d-7}(h),
\qquad
\widehat Y_d^{OFF}(h).
\]

Separate weekly score autoregression:
\[
\xi_{dk}=\alpha_k+\beta_k\xi_{d-7,k}+u_{dk},
\qquad k=1,\ldots,K.
\]

Weekly FPCA VAR:
\[
\xi_d=a+A_7\xi_{d-7}+u_d,
\qquad A_7\in{\cal R}^{K\times K}.
\]

FPCA VAR(1,7):
\[
\xi_d=a+A_1\xi_{d-1}+A_7\xi_{d-7}+u_d,
\qquad A_1,A_7\in{\cal R}^{K\times K}.
\]

Augmented score model:
\[
\xi_d=a+A_1\xi_{d-1}+A_7\xi_{d-7}+Bz_d+u_d,
\]
where \(z_d\) is predetermined at date \(d\).

Residual-extension forecast:
\[
\widehat Y_d^{res}(h)
=
\widehat\mu(h)+\widehat{\Phi}(h)'\widehat\xi_d+\lambda_d\widehat r_d(h).
\]

\section*{Predictor construction}

The largest predictor vector is
\[
z_d=
\left(
z_d^{cal'},z_d^{hol'},z_d^{shape'},z_d^{deriv'}
\right)'.
\]

Ramp and early-load variables:
\[
R^m_{d-1}=Y_{d-1}(9)-Y_{d-1}(6),
\qquad
R^m_{d-7}=Y_{d-7}(9)-Y_{d-7}(6),
\]
\[
L^e_{d-1}=\frac{1}{3}\sum_{h\in\{5,6,7\}}Y_{d-1}(h).
\]

Calendar controls:
\[
z_d^{cal}
=
\left(
weekend_d,
\sin\frac{2\pi doy_d}{365.25},
\cos\frac{2\pi doy_d}{365.25}
\right)'.
\]

Holiday controls include fixed holidays, religious holidays, holiday eve, post-holiday, near-holiday, and bridge-day indicators:
\[
Bridge_d=
1\{dow_d\in\{Mon,Fri\}\}
1\{d\notin{\cal H}\}
1\{d-1\in{\cal H}\ {\rm or}\ d+1\in{\cal H}\}.
\]

Lagged load-shape variables:
\[
\bar Y_{d-1}=\frac{1}{24}\sum_{h=0}^{23}Y_{d-1}(h),
\qquad
P_{d-1}=\max_hY_{d-1}(h),
\qquad
M_{d-1}=\min_hY_{d-1}(h),
\]
\[
G_{d-1}=P_{d-1}-M_{d-1},
\qquad
\bar Y_d^{(q)}
=
\frac{1}{q}\sum_{j=1}^{q}
\left(
\frac{1}{24}\sum_{h=0}^{23}Y_{d-j}(h)
\right),
\quad q\in\{7,28\}.
\]

Derivative features are computed from the smoothed curve,
\[
DY_d(h)=\frac{d}{dh}Y_d(h),
\]
and summarized by
\[
\overline{DY}_{d-1},\quad
SD(DY_{d-1}),\quad
\max_hDY_{d-1}(h),\quad
\min_hDY_{d-1}(h),\quad
\max_hDY_{d-1}(h)-\min_hDY_{d-1}(h).
\]

Targeted intraday changes:
\[
R^e_{d-1}=Y_{d-1}(7)-Y_{d-1}(5),
\qquad
R^{ev}_{d-1}=Y_{d-1}(20)-Y_{d-1}(17),
\qquad
R^{end}_{d-1}=Y_{d-1}(23)-Y_{d-1}(20).
\]

\section*{Evaluation}

For model \(m\),
\[
e_{dh}^{(m)}=Y_d(h)-\widehat Y_d^{(m)}(h).
\]

\[
MAE_m=\frac{1}{N}\sum_{d,h}|e_{dh}^{(m)}|,
\qquad
RMSE_m=
\left[
\frac{1}{N}\sum_{d,h}(e_{dh}^{(m)})^2
\right]^{1/2},
\qquad
Bias_m=\frac{1}{N}\sum_{d,h}e_{dh}^{(m)}.
\]

\[
Gain_m^{MAE}=100\frac{MAE_{OFF}-MAE_m}{MAE_{OFF}},
\qquad
Gain_m^{RMSE}=100\frac{RMSE_{OFF}-RMSE_m}{RMSE_{OFF}}.
\]

\section*{Preliminary results}

The rows below should be read as a model-building path. The first two rows are
benchmarks. Each later row adds one modeling idea: FPCA score dynamics, calendar and
ramp information, holiday/load-shape information, and derivative information.

\begin{center}
\fontsize{8.2}{9.5}\selectfont
\begin{tabular}{p{0.41\textwidth}rrrrr}
\hline
Model & MAE & RMSE & Bias & MAE gain & RMSE gain\\
\hline
Seasonal naive \(Y_{d-7}\) & 1912.46 & 3129.98 & 39.85 & -57.60\% & -79.85\%\\
Official forecast & 1213.50 & 1740.30 & 466.87 & 0.00\% & 0.00\%\\
FPCA VAR(1,7) scores & 1442.61 & 2018.30 & 160.13 & -18.88\% & -15.97\%\\
FPCA VAR(1,7) + ramp/calendar & 1238.37 & 1716.30 & 110.12 & -2.05\% & 1.38\%\\
FPCA VAR(1,7) + holiday/load-shape & 1000.62 & 1353.56 & -54.71 & 17.54\% & 22.22\%\\
FPCA VAR(1,7) + holiday/load-shape/derivative & {\bf 980.38} & {\bf 1324.40} & {\bf 102.56} & {\bf 19.21\%} & {\bf 23.90\%}\\
\hline
\end{tabular}
\end{center}

\fontsize{9.2}{10.8}\selectfont

The final model improves on the official forecast by \(19.21\%\) in MAE and \(23.90\%\) in RMSE. Relative to the weekly seasonal naive benchmark,
\[
100\frac{1912.46-980.38}{1912.46}=48.74\%.
\]

\section*{Discussion: K and residual correction}

The main report uses \(K=2\), because this follows the 95\% cumulative-variance rule
for the load FPCA and gives the most parsimonious score model. The cost of this
choice is some loss of accuracy relative to \(K=4\). The \(K=4\) no-residual model
has lower overall MAE and RMSE, and it beats the official forecast in 23 of 24 hours
by MAE and 24 of 24 hours by RMSE. The \(K=2\) no-residual model wins in 19 of 24
hours by MAE and 21 of 24 hours by RMSE.

\begin{center}
\fontsize{7.6}{8.8}\selectfont
\begin{tabular}{p{0.26\textwidth}rrrrrr}
\hline
Specification & \(K\) & Resid. step & MAE & RMSE & MAE gain & RMSE gain\\
\hline
Holiday/load/derivative & 2 & No & 980.38 & 1324.40 & 19.21\% & 23.90\%\\
Holiday/load/derivative & 2 & Yes & 652.90 & 982.60 & 46.20\% & 43.54\%\\
Holiday/load/derivative & 4 & No & 800.99 & 1109.69 & 33.99\% & 36.24\%\\
Holiday/load/derivative & 4 & Yes & {\bf 640.10} & {\bf 967.82} & {\bf 47.25\%} & {\bf 44.39\%}\\
\hline
\end{tabular}
\end{center}

\fontsize{9.2}{10.8}\selectfont

The table suggests a tradeoff between parsimony and predictive accuracy.
The \(K=2\) specification is the conservative baseline, while \(K=4\) retains
additional intraday shape variation and improves the no-residual forecast.
I therefore keep \(K=2\) as the main specification, but report \(K=4\) as a robustness check.

The second-stage residual correction should be interpreted even more cautiously, especially its methodological appropriateness.
The closest related papers are not functional forecasting papers,
but papers on residual post-processing and error compensation in electricity forecasting.
\citet{xie_hong_2016} use forecast combination and residual simulation in probabilistic electricity load forecasting.
\citet{wang_etal_2012} use residual modification models to improve seasonal ARIMA electricity demand forecasts.
\citet{kontogiannis_etal_2022} estimate hourly residual error sequences to refine day-ahead electricity price forecasts.
\citet{ghimire_etal_2024} also use a two-step forecasting design in which an initial electricity price forecast is improved by an error-compensation model.
These papers motivate the idea of learning from past forecast errors, but I do not yet claim that they validate the exact functional residual correction used here.
For this reason, I treat the residual correction as a promising empirical extension rather than as the main model.

For target day
\(d\), the base model is first estimated exactly as before. Then fitted residual
curves are computed only for days strictly before the target day,
\[
\widehat e_s(h)=Y_s(h)-\widehat Y_s^{base}(h),
\qquad s<d.
\]
The residual second step estimates a correction curve \(\widehat r_d(h)\) from these
past residuals and applies
\[
\widehat Y_d^{res}(h)=\widehat Y_d^{base}(h)+\lambda_d\widehat r_d(h).
\]
Both the residual regression and the validation choice of \(\lambda_d\) use only
information dated strictly before \(d\). Thus the extension is still rolling-origin
and is designed to avoid information leakage.



\begin{center}
\fontsize{8.2}{9.5}\selectfont
\begin{tabular}{p{0.43\textwidth}rrrrr}
\hline
Model & MAE & RMSE & Bias & MAE gain & RMSE gain\\
\hline
Best base model & 980.38 & 1324.40 & 102.56 & 19.21\% & 23.90\%\\
Base model + residual second step & {\bf 652.90} & {\bf 982.60} & {\bf 90.03} & {\bf 46.20\%} & {\bf 43.54\%}\\
\hline
\end{tabular}
\end{center}

\fontsize{9.2}{10.8}\selectfont


\section*{Conclusion}

This preliminary report treats Turkish hourly electricity load as daily curves and
evaluates whether an FPCA-based rolling-origin model can improve on the official
day-ahead forecast. The conservative \(K=2\) specification improves on the official
forecast by \(19.21\%\) in MAE and \(23.90\%\) in RMSE.

The results suggest that intraday load-shape information, calendar effects, holidays,
and derivative/ramp features contain useful predictive information. However, the
model specification is still preliminary. Further work is needed to decide how many
components should be retained, whether the residual correction is methodologically
appropriate, and how the forecast gains can be connected to the economic intuition
behind intraday demand patterns and system operation.



\begin{thebibliography}{99}

\bibitem[Xie and Hong(2016)]{xie_hong_2016}
Xie, J. and Hong, T. (2016).
GEFCom2014 probabilistic electric load forecasting: An integrated solution with
forecast combination and residual simulation.
\textit{International Journal of Forecasting}, 32(3), 1012--1016.
\href{https://doi.org/10.1016/j.ijforecast.2015.11.005}{doi:10.1016/j.ijforecast.2015.11.005}.

\bibitem[Wang et al.(2012)]{wang_etal_2012}
Wang, Y., Wang, J., Zhao, G. and Dong, Y. (2012).
Application of residual modification approach in seasonal ARIMA for electricity
demand forecasting: A case study of China.
\textit{Energy Policy}, 48, 284--294.
\href{https://doi.org/10.1016/j.enpol.2012.05.026}{doi:10.1016/j.enpol.2012.05.026}.

\bibitem[Kontogiannis et al.(2022)]{kontogiannis_etal_2022}
Kontogiannis, D., Bargiotas, D., Daskalopulu, A., Arvanitidis, A. I. and
Tsoukalas, L. H. (2022).
Error compensation enhanced day-ahead electricity price forecasting.
\textit{Energies}, 15(4), 1466.
\href{https://doi.org/10.3390/en15041466}{doi:10.3390/en15041466}.

\bibitem[Ghimire et al.(2024)]{ghimire_etal_2024}
Ghimire, S., Deo, R. C., Casillas-Pérez, D. and Salcedo-Sanz, S. (2024).
Two-step deep learning framework with error compensation technique for short-term,
half-hourly electricity price forecasting.
\textit{Applied Energy}, 353, 122059.
\href{https://doi.org/10.1016/j.apenergy.2023.122059}{doi:10.1016/j.apenergy.2023.122059}.

\bibitem[Kokoszka and Reimherr(2017)]{kokoszka_reimherr_2017}
Kokoszka, P. and Reimherr, M. (2017).
\textit{Introduction to Functional Data Analysis}.
CRC Press.

\end{thebibliography}






\end{document}
